% =====================================================================
% EdgeFlow 用户手册 · 第 6 章 设备接入
% 事实依据：docs/API-SPEC.md、docs/MAPPER-GUIDE.md、docs/MODBUS-GUIDE.md、
%           docs/EVENTBUS-GUIDE.md、edge/pkg/devicetwin、cmd/edgecore
% =====================================================================
\chapter{设备接入}
\label{ch:6}

EdgeFlow 通过\textbf{设备影子}（DeviceTwin）统一管理边缘设备：设备数据在
边缘汇聚、周期上报云端；云端指令经边缘下发到设备。本章说明设备模型、
指令链路、Mapper 框架（设备映射器）与数据面。

\section{设备模型：影子}

每台设备在边缘维护一份\textbf{影子}（Shadow / DeviceTwin），包含两类属性：

\begin{description}
  \item[\texttt{properties}（实际值）] 设备实际上报的属性值（如温度、湿度），
        由 Mapper 采样写入，周期上报云端；
  \item[\texttt{desired}（期望值）] 云端期望设备达到的属性值，通过指令下发
        写入，云端与边缘各存一份。
\end{description}

影子是边缘侧设备数据的\textbf{汇聚点}：无论设备通过何种协议接入（MQTT、
Modbus\ldots），云端看到的都是统一的影子视图。边缘默认每
\texttt{EDGEFLOW\_EDGECORE\_DEVICE\_REPORT\_INTERVAL}（30 秒）把影子快照
上报云端。

\section{设备指令下发链路}

云端下发指令走\textbf{可靠投递}通道（见第~\ref{ch:5}~章），完整链路为：

\begin{enumerate}
  \item 调用方 \texttt{POST /api/v1/nodes/\{nodeID\}/device-command}，
        携带 \texttt{deviceName} 与 \texttt{property}（期望值）；
  \item 云端经 CloudHub 将 DeviceCommand 消息可靠送达边缘（Ack 确认后
        API 返回 \texttt{200}）；
  \item 边缘按 \texttt{deviceName} 路由到对应 Mapper 直接执行指令（MQTT 模式下另订阅本地指令主题 \texttt{devices/<ns>/<device>/command}，支持数据面直接下发），设备属性发生变化；
  \item 边缘影子 \texttt{properties} 更新，下一上报周期同步到云端。
\end{enumerate}

\begin{lstlisting}
# 下发设备指令：把 sensor-01 的目标温度设为 25
curl -X POST http://127.0.0.1:8080/api/v1/nodes/edge-node-1/device-command \
  -H 'Content-Type: application/json' \
  -d '{"deviceName":"sensor-01","property":"targetTemp","value":25}'
# 期望响应：{"status":"ok","acked":true}
\end{lstlisting}

常见失败语义：节点离线 → \texttt{404}（未送达，无需重试）；边缘拒绝 →
\texttt{502}（重试无意义）；确认超时 → \texttt{504}（可重试，边缘幂等去重）。
详见附录~B。

\section{Mapper 框架}

\textbf{Mapper}（设备映射器）负责把具体设备的协议翻译成 EdgeFlow 的统一
设备模型。Mapper 注册到边缘的 \texttt{MapperRegistry}，按
\texttt{deviceName} 路由设备消息。v0.1.0 内置两个 Mapper（\texttt{modbus} 需显式启用，见下）：

\subsection{Mapper 装配开关（v0.2.0）}

v0.2.0 起，Mapper 装配由开关 \texttt{EDGEFLOW\_EDGECORE\_ENABLE\_MAPPER}
控制，默认 \texttt{true}（与旧版本行为一致）：

\begin{itemize}
  \item \texttt{true}（或 \texttt{1}/\texttt{on}/\texttt{yes}，大小写不敏感）：装配全部 Mapper，正常采集与指令执行；
  \item \texttt{false}（或 \texttt{0}/\texttt{off}/\texttt{no}）：\textbf{不注册任何 Mapper}、不启动采集循环；云端下发的设备指令\textbf{仅更新影子 \texttt{Twin.Desired}，不下发到设备}——即\textbf{纯影子模式}。
\end{itemize}

\subsection{mock\_sensor：模拟传感器}

内置模拟温湿度传感器，用于验证全链路，无需真实硬件：

\begin{itemize}
  \item 设备名：\texttt{sensor-01}；
  \item 采集周期：每 \texttt{2 秒}波动一次温度/湿度；
  \item 数据面：MQTT（遥测主题 \texttt{devices/default/sensor-01/telemetry}）；
  \item 支持指令：如 \texttt{targetTemp}（目标温度）写入期望值。
\end{itemize}

\subsection{Modbus TCP Mapper}

通过 \texttt{goburrow/modbus} 客户端库（社区标准 Modbus 客户端）对接
Modbus TCP 设备：

\begin{itemize}
  \item 设备名：\texttt{mb-sensor-01}（unit ID 1），采集温度/湿度保持寄存器；
  \item 自带 \texttt{modbussim} 模拟器，无需真实设备即可联调；
  \item \textbf{操作台账}（op-ledger）：每次读写操作落 SQLite 台账，可查
        设备 ID、方向（\texttt{up}/\texttt{down}）与时间，便于审计排障；
  \item 指令链路与 mock\_sensor 一致：云端下发
        \texttt{DeviceCommand\{deviceName:"mb-sensor-01",\ldots\}}，
        Mapper 按指令写寄存器。
\end{itemize}

\textbf{启用方式}：Modbus Mapper 为可选接入，必须显式设置环境变量
\texttt{EDGEFLOW\_MODBUS\_ADDR}（如 \texttt{127.0.0.1:15020} 指向本机模拟器，
或 \texttt{<设备IP>:502} 指向真实 Modbus 设备）后启动 edgecore，该 Mapper
才会注册；未设置时 \texttt{mb-sensor-01} 不存在。

\subsection{设备命名空间（v0.2.0）}

v0.2.0 起，设备按\textbf{命名空间}（namespace）分组，路由键为
\texttt{namespace/deviceName}，同名设备可在不同命名空间共存：

\begin{itemize}
  \item \textbf{解析优先级}（三级）：Mapper 的 \texttt{WithNamespace} 选项 >
        环境变量 \texttt{EDGEFLOW\_MODBUS\_NAMESPACE} > 默认 \texttt{default}；
  \item \textbf{注册表隔离}：MapperRegistry 按命名空间隔离路由，同名设备
        分属不同命名空间互不干扰；
  \item \textbf{指令路由}：\texttt{device-command} 请求的 \texttt{namespace}
        参与路由，命名空间不匹配任何 Mapper 时边缘拒绝 \texttt{502}；
  \item \textbf{采集汇入}（v0.3.0 修复）：采集数据按\texttt{mapper 自身
        命名空间}汇入影子，多命名空间部署不会出现 \texttt{default/} 双条目
        错位。
\end{itemize}

\subsection{OPC-UA 接入状态（v0.3.0 起）}

v0.3.0 交付\textbf{OPC-UA 第一阶段}：UA Binary 协议栈核心（\texttt{pkg/opcua}，
零第三方依赖，纯标准库实现）。当前能力与边界如下，接入真实 OPC-UA 设备
排后续版本：

\begin{itemize}
  \item \textbf{已交付}：UA Binary 编解码（NodeId 全编码形式、Variant、
        DataValue、DateTime 等类型）与 \texttt{HEL}/\texttt{ACK} 握手
        （\texttt{Dial}/\texttt{DialTimeout}/\texttt{ReadMessage}/
        \texttt{WriteMessage}）；
  \item \textbf{未实现}：设备读写服务（Read/Write/Subscribe）、
        \texttt{SecureChannel}（当前为裸传输，ChannelId=0）、节点模型、
        Discovery、Sign/SignAndEncrypt 安全策略——因此\textbf{尚无 OPC-UA
        Mapper 接入}，设备接入链路未打通；
  \item \textbf{安全边界}：仅支持 \texttt{SecurityPolicy None}（\textbf{明文}），
        \textbf{严禁暴露到不可信网络}，仅限可信隔离网络使用；
  \item \textbf{验证状态}：未与 open62541 / node-opcua 等第三方栈做过互
        操作验证。
\end{itemize}

\subsection{Mapper 的两种工作模式}

Mapper 由 edgecore 装配时统一注册，按 EventBus 是否可用分为两种模式：

\begin{description}
  \item[MQTT 模式] EventBus 连接成功：遥测发布到 MQTT 主题、订阅设备指令
        主题，数据面完整（broker 掉线后 paho 自动重连、订阅自动恢复）；
  \item[纯本地模式] MQTT 不可用（未装配/连接失败降级）：采集照常写入
        影子、本地指令照常执行，仅 MQTT 数据面不可用；云边通道
        （DeviceCommand/DeviceReport）不受影响。
\end{description}

两种模式对云端完全透明：云端只看到影子视图，不感知边缘内部数据面
形态。这也是第~\ref{ch:9}~章「EventBus 连接失败不必停机」的机制基础。

\section{EventBus 数据面与降级}

边缘内部通过 \textbf{EventBus}（MQTT，本机 1883）承载设备数据面：
设备/Mapper 之间的遥测与指令走 MQTT 主题，云边管理消息仍走 WebSocket
通道，二者互不干扰。

\begin{table}[htbp]
\centering
\caption{EventBus 主题约定（QoS 1）}
\label{tab:6-topics}
\begin{tabularx}{\textwidth}{l l X}
\toprule
\textbf{主题} & \textbf{方向} & \textbf{用途} \\
\midrule
\texttt{devices/<ns>/<device>/telemetry} & 设备 → 边缘 & 设备遥测上报 \\
\texttt{devices/<ns>/<device>/command} & 边缘 → 设备 & 指令下发 \\
\texttt{edgeflow/<module>/<event>} & 模块间 & 内部事件（预留） \\
\bottomrule
\end{tabularx}
\end{table}

\textbf{降级行为}：MQTT broker 不可用或连接超时（默认 5 秒）时，Mapper
自动降级\textbf{纯本地模式}——设备采集与本地指令照常工作，仅数据面
MQTT 发布/订阅不可用；云边通道（DeviceCommand/DeviceReport）不受影响。
已建立连接后的掉线会自动重连（QoS 1，订阅自动恢复）；若为\textbf{启动时}连接失败而降级，则降级是装配期决策，broker 恢复后需重启 edgecore 才能启用 MQTT 数据面。

\begin{lstlisting}
# 订阅查看 sensor-01 遥测（本机 mosquitto broker）
mosquitto_sub -p 1883 -t 'devices/default/sensor-01/telemetry' -v
# 通过 MQTT 数据面直接下发指令
mosquitto_pub -p 1883 -t 'devices/default/sensor-01/command' \
  -m '{"deviceName":"sensor-01","property":"targetTemp","value":23}'
\end{lstlisting}

\section{设备状态查询}

云端提供统一查询入口，返回设备的 \texttt{properties} 与 \texttt{desired}：

\begin{lstlisting}
# 全部设备状态
curl http://127.0.0.1:8080/api/v1/devices
# 单节点设备状态
curl http://127.0.0.1:8080/api/v1/nodes/{nodeID}/devices
\end{lstlisting}

\begin{itemize}
  \item 节点不存在 → \texttt{404}；
  \item 节点存在但无设备 → \texttt{200} + 空 \texttt{items}（可无分支遍历）；
  \item 开启 Token 认证时需携带 \texttt{Authorization: Bearer <token>}
        请求头（见第~\ref{ch:8}~章）。
\end{itemize}

\textbf{持久化注（v0.4.0 起）}：设备 \texttt{desired}（期望值）随云端
etcd 写穿\textbf{持久化}，跨云端重启保留；上报属性（\texttt{properties}）
在云端重启后短暂清空，由边缘\textbf{不超过 1 个上报周期}内自动补报自愈。

\section{设备接入检查清单}

接入一台新设备时，按以下顺序核对：

\begin{enumerate}
  \item \textbf{数据面}：设备/Mapper 与 broker（或本地模式）连通，遥测
        主题有数据（\texttt{mosquitto\_sub} 可观测）；
  \item \textbf{影子}：edgecore 日志出现 Mapper 注册成功，影子
        \texttt{properties} 有值；
  \item \textbf{上报}：等一个上报周期（默认 30s），云端
        \texttt{GET /api/v1/devices} 可见设备；
  \item \textbf{指令}：下发一条 \texttt{device-command}，确认 200 + 设备
        属性变化 + 期望值写入云端。
\end{enumerate}
